\documentclass[aspectratio=169]{beamer}

% ── Theme & Colors ────────────────────────────────────────────────────────────
\usetheme{Madrid}
\usecolortheme{default}

\definecolor{utepblue}{RGB}{0, 56, 101}      % UTEP navy
\definecolor{utepgold}{RGB}{255, 200, 37}     % UTEP gold
\definecolor{lightgray}{RGB}{245, 245, 245}

\setbeamercolor{palette primary}{bg=utepblue, fg=white}
\setbeamercolor{palette secondary}{bg=utepblue!80, fg=white}
\setbeamercolor{palette tertiary}{bg=utepblue!60, fg=white}
\setbeamercolor{palette quaternary}{bg=utepblue, fg=white}
\setbeamercolor{structure}{fg=utepblue}
\setbeamercolor{section in toc}{fg=utepblue}
\setbeamercolor{frametitle}{bg=utepblue, fg=white}
\setbeamercolor{title}{bg=utepblue, fg=white}
\setbeamercolor{block title}{bg=utepblue, fg=white}
\setbeamercolor{block body}{bg=lightgray}
\setbeamercolor{alerted text}{fg=utepgold!80!black}

\setbeamertemplate{navigation symbols}{}
\setbeamertemplate{footline}[frame number]
\setbeamertemplate{itemize items}[circle]

% ── Packages ──────────────────────────────────────────────────────────────────
\usepackage{booktabs}
\usepackage{graphicx}
\usepackage{amsmath}
\usepackage{multirow}
\usepackage{xcolor}
\usepackage{tikz}

% ── Title Information ─────────────────────────────────────────────────────────
\title[Twitter Algorithmic Bias]{Algorithmic Amplification and Content Bias on Twitter/X:\\
\small Evidence from the Musk Acquisition}
\subtitle{Panel Data and Discrete Choice Models --- Spring 2026}
\author{Alex Gamez}
\institute{The University of Texas at El Paso}
\date{April 2026}

% ─────────────────────────────────────────────────────────────────────────────
\begin{document}

% ── SLIDE 1 — Title ───────────────────────────────────────────────────────────
\begin{frame}
\titlepage
\end{frame}

% ── SLIDE 2 — Motivation ──────────────────────────────────────────────────────
\begin{frame}{Motivation}
\textbf{Social media algorithms shape what people see --- and what they believe.}
\vspace{0.5em}
\begin{itemize}
    \item Algorithmic feed on X (Twitter) for 7 weeks shifted users' political opinions toward more conservative views \hfill \textcolor{gray}{\small Gauthier et al., \textit{Nature} (2026)}
    \vspace{0.3em}
    \item Facebook's ``I Voted'' message increased voter turnout by 0.39pp ($\approx$340,000 additional votes) in the 2010 U.S. election \hfill \textcolor{gray}{\small Bond et al., \textit{Nature} (2012)}
    \vspace{0.3em}
    \item Facebook's engagement-based algorithm amplified hate speech against the Rohingya, contributing to mass violence and 700,000+ displaced people \hfill \textcolor{gray}{\small Amnesty International (2022)}
    \vspace{0.3em}
    \item TikTok's algorithm increased exposure to misogynistic content by 4$\times$ in five days for teen accounts \hfill \textcolor{gray}{\small UCL \& Univ.\ of Kent (2024)}
    \vspace{0.3em}
    \item Researchers document misogynistic ideologies ``moving off screens and into schools,'' shaping real-world behavior among boys \hfill \textcolor{gray}{\small UCL (2024)}
\end{itemize}
\end{frame}

% ── SLIDE 3 — Research Question ───────────────────────────────────────────────
\begin{frame}{Research Question}

\begin{block}{Central Question}
Did Elon Musk's acquisition of Twitter on October 27, 2022 \textbf{causally change} the content categories amplified in Twitter's trending algorithm --- relative to what organic user interest would have produced?
\end{block}

\vspace{0.8em}

\begin{columns}[t]
    \column{0.33\textwidth}
    \centering
    \textbf{Treatment}\\[0.3em]
    Musk completes acquisition\\
    \textbf{October 27, 2022}\\[0.3em]
    {\small Mass layoffs, policy reversals, and algorithmic changes followed immediately}

    \column{0.33\textwidth}
    \centering
    \textbf{Outcome}\\[0.3em]
    $\log(\text{daily share of trending topics in category } c)$\\[0.3em]
    {\small Measured for 16 content categories across a 4-year window}

    \column{0.33\textwidth}
    \centering
    \textbf{Counterfactual}\\[0.3em]
    $\log(\text{Reddit posts in matched subreddit})$\\[0.3em]
    {\small Proxies organic user interest unaffected by the ownership change}
\end{columns}

\vspace{0.5em}
\small \textit{Identifying assumption:} absent the acquisition, Twitter trending content would have tracked Reddit's organic interest trajectory --- the parallel trends assumption.
\end{frame}

% ── SLIDE 4 — Data & Sources ──────────────────────────────────────────────────
\begin{frame}{Data \& Sources}
\begin{columns}[t]
    \column{0.33\textwidth}
    \textbf{\textcolor{utepblue}{Twitter Trending}}
    \begin{itemize}
        \item 4-year panel\\Oct 2020 – Oct 2024
        \item $\approx$695,000 hourly observations
        \item 18 content categories via NLP lexicon
    \end{itemize}

    \column{0.33\textwidth}
    \textbf{\textcolor{utepblue}{Reddit (Control)}}
    \begin{itemize}
        \item Arctic Shift time series API
        \item Daily post counts for 25+ matched subreddits
        \item Full 4-year date range
        \item 35,900+ subreddit-day observations
    \end{itemize}

    \column{0.33\textwidth}
    \textbf{\textcolor{utepblue}{Content Classification}}
    \begin{itemize}
        \item CamelCase-aware keyword lexicon
        \item 18 categories
        \item $\approx$38\% of unique topics classified; covers majority of high-frequency trending slots
    \end{itemize}
\end{columns}
\end{frame}

% ── SLIDE 5 — Data Collection Challenges ─────────────────────────────────────
\begin{frame}{Data Collection: Challenges \& Solution}
\begin{columns}[t]
    \column{0.48\textwidth}
    \textbf{\textcolor{red!70!black}{First Approach (Problems)}}
    \begin{itemize}
        \item Attempted large samples by week/month, classified into conservative / neutral / progressive using online dictionary lexicons
        \item Twitter API \textbf{could not provide} historical trending data at the required scale
        \item Severe data limitations prevented longitudinal analysis
    \end{itemize}

    \column{0.48\textwidth}
    \textbf{\textcolor{utepblue}{Solution}}
    \begin{itemize}
        \item Shifted to archived Twitter trending data with content-category classification
        \item Arctic Shift API --- full 4-year range in 2 calls per subreddit
        \item No post cap --- accurate counts across all subreddits
        \item Outcome reframed: \textit{category share shift} rather than ideological labeling
    \end{itemize}
\end{columns}
\end{frame}

% ── SLIDE 6 — Identification Strategy ────────────────────────────────────────
\begin{frame}{Identification Strategy: Difference-in-Differences}

\begin{block}{Log-Linear DiD Model (HC3 robust standard errors)}
\[
\log(y_{it}) = \beta_0 + \beta_1 \cdot \text{Twitter}_i + \beta_2 \cdot \text{Post}_t + \underbrace{\beta_3 \cdot (\text{Twitter}_i \times \text{Post}_t)}_{\text{causal effect}} + \varepsilon_{it}
\]
\end{block}

\vspace{0.4em}
\begin{columns}[t]
    \column{0.55\textwidth}
    \begin{itemize}
        \item $\log(y_{it})$ = log(Twitter topic share $+ \epsilon$) \textbf{or} log(Reddit posts $+ 1$)
        \item $\text{Twitter}_i = 1$ for Twitter, 0 for Reddit (two units stacked per category)
        \item $\text{Post}_t = 1$ after October 27, 2022
        \item $\hat{\beta}_3$ = log-point DiD $\approx$ \% effect for small values
        \item HC3 heteroskedasticity-robust standard errors
    \end{itemize}

    \column{0.42\textwidth}
    \begin{alertblock}{Interpretation}
    \textbf{Positive} $\hat{\beta}_3$: Twitter amplified beyond organic interest\\[0.4em]
    \textbf{Negative} $\hat{\beta}_3$: Twitter suppressed below organic interest
    \end{alertblock}
\end{columns}
\end{frame}

% ── SLIDE 6b — Identification Assumptions ────────────────────────────────────
\begin{frame}{Identification Assumptions}

\begin{enumerate}
    \setlength\itemsep{0.6em}

    \item \textbf{Parallel Trends}\\
    In the absence of the Musk acquisition, Twitter and Reddit trending content would have followed the same trajectory. \textit{Tested visually (pre-period trends) and via event study pre-quarters $\approx 0$.}

    \item \textbf{No Anticipation}\\
    Platform behavior did not change before October 27, 2022 in response to the anticipated acquisition. \textit{Tested: no structural break detected in pre-period.}

    \item \textbf{Stable Unit Treatment Value (SUTVA)}\\
    Reddit's activity is not itself affected by Twitter's algorithmic changes. Platforms operate independently; Reddit has no comparable ownership shock in this window.

    \item \textbf{Exogeneity of Treatment Timing}\\
    The October 27, 2022 date is determined by court proceedings and financing --- not by content trends. Musk was compelled to close by a Delaware court order after attempting to withdraw.

    \item \textbf{Common Support}\\
    Both Twitter and Reddit have non-zero activity in all categories throughout the 4-year window.

\end{enumerate}
\end{frame}

% ── SLIDE 7 — Parallel Trends & Validity ─────────────────────────────────────
\begin{frame}{Parallel Trends \& Validity}

\begin{columns}[t]
    \column{0.55\textwidth}
    % Upload out/figures/parallel_trends/wrestling.png (or any representative category)
    \includegraphics[width=\textwidth]{figures/parallel_trends/wrestling.png}
    \vspace{0.2em}
    \includegraphics[width=\textwidth]{figures/parallel_trends/sports_womens.png}

    \column{0.42\textwidth}
    \textbf{What to look for}
    \begin{itemize}
        \item Pre-treatment (left of dashed line): Twitter and Reddit lines move \textbf{together} $\Rightarrow$ parallel trends satisfied
        \item Post-treatment (right of dashed line): lines \textbf{diverge} $\Rightarrow$ evidence of algorithmic shift
    \end{itemize}
    \vspace{0.5em}
    \textbf{Validity checks}
    \begin{itemize}
        \item Placebo test: false treatment dates yield null results
        \item Anticipation test: no structural break before Oct 27, 2022
        \item Event study: pre-period quarters cluster near zero
    \end{itemize}
    \vspace{0.3em}
    {\small \textit{Y-axis:} log outcome, demeaned to pre-period mean so both series start at 0.}
\end{columns}
\end{frame}

% ── SLIDE 8 — Content Classification ─────────────────────────────────────────
\begin{frame}{Content Classification}

\begin{columns}[t]
    \column{0.48\textwidth}
    \textbf{18 Categories}
    \begin{itemize}\small
        \item Wrestling, Combat Sports
        \item NBA, NFL, MLB, NHL, Soccer, College Sports, Women's Sports
        \item Entertainment, Reality TV, Taylor Swift, Fandom
        \item Tech \& Gaming
        \item News \& Events, Politics
        \item Holidays, Social Filler
    \end{itemize}

    \column{0.48\textwidth}
    % INSERT FIGURE: out/figures/category_shift.png
    \fbox{\parbox{\textwidth}{\centering\vspace{3.5cm}
    \texttt{[INSERT FIGURE: category\_shift.png]}
    \vspace{3.5cm}}}
\end{columns}

\vspace{0.3em}
\small CamelCase-aware tokenizer: \texttt{\#WWERaw} $\to$ \{wwe, raw\} $\to$ ``Wrestling''.
Matched Reddit subreddit control per category (e.g., r/SquaredCircle for Wrestling).
\end{frame}

% ── SLIDE 9 — Main Results ────────────────────────────────────────────────────
\begin{frame}{Main Results}

\begin{center}
% INSERT FIGURE: out/figures/category_did.png
\fbox{\parbox{0.92\textwidth}{\centering\vspace{4.5cm}
\texttt{[INSERT FIGURE: category\_did.png]}\\
$\hat{\beta}_3$ estimates with 95\% CI. HC3 robust standard errors.\\
Positive = amplified relative to Reddit baseline. Negative = suppressed.
\vspace{4.5cm}}}
\end{center}
\end{frame}

% ── SLIDE 10 — Interpreting the Results ──────────────────────────────────────
\begin{frame}{Interpreting the Results}
\begin{columns}[t]
    \column{0.48\textwidth}
    \textbf{\textcolor{utepblue}{Amplified} (positive $\hat{\beta}_3$)}
    \begin{itemize}
        \item \textbf{Wrestling} $+63$pp*** --- WWE trending far above Reddit's organic growth
        \item \textbf{News \& Politics} $+53$pp*** --- political content boosted post-acquisition
        \item \textbf{Tech \& Gaming} $+39$pp*** --- consistent with Musk's tech-focused audience
        \item \textbf{Combat Sports} $+49$pp*** --- UFC/boxing trending accelerated
    \end{itemize}

    \column{0.48\textwidth}
    \textbf{\textcolor{red!70!black}{Suppressed} (negative $\hat{\beta}_3$)}
    \begin{itemize}
        \item \textbf{Taylor Swift} $-106$pp*** --- suppressed despite massive real-world growth (Eras Tour)
        \item \textbf{Women's Sports} $-77$pp*** --- WNBA/NWSL Reddit surged (Caitlin Clark); Twitter did not follow
        \item \textbf{Fandom} $-20$pp*** --- K-pop/Anime organically declined; Twitter followed
    \end{itemize}
\end{columns}
\end{frame}

% ── SLIDE 11 — Robustness Checks ─────────────────────────────────────────────
\begin{frame}{Robustness Checks}

\begin{columns}[t]
    \column{0.48\textwidth}
    % INSERT FIGURE: out/figures/did_multiwindow.png (if generated)
    \fbox{\parbox{\textwidth}{\centering\vspace{3.5cm}
    \texttt{[INSERT FIGURE: multiwindow results]}
    \vspace{3.5cm}}}

    \column{0.48\textwidth}
    \textbf{Multi-Window DiD}
    \begin{itemize}
        \item 6-month window: $\hat{\beta}_3 = +5.5$pp (n.s.)
        \item 1-year window: $\hat{\beta}_3 = +6.9$pp (*)
        \item 2-year window: $\hat{\beta}_3 = +2.9$pp (n.s.)
    \end{itemize}

    \vspace{0.5em}
    \textbf{Alternative Measures}
    \begin{itemize}
        \item Volume-weighted Twitter share: consistent
        \item Score-weighted Reddit: consistent
        \item Placebo tests: null results at false treatment dates
    \end{itemize}
\end{columns}
\end{frame}

% ── SLIDE 12 — Limitations ───────────────────────────────────────────────────
\begin{frame}{Limitations}
\begin{itemize}
    \item \textbf{Reddit $\neq$ perfect counterfactual} --- platform demographics differ from Twitter; may not fully proxy Twitter's organic trajectory
    \vspace{0.4em}
    \item \textbf{$\approx$62\% of topics unclassified} --- long-tail hashtags create potential selection bias toward high-salience trending events
    \vspace{0.4em}
    \item \textbf{Confounds:} Caitlin Clark era (Women's Sports) and Eras Tour (Taylor Swift) coincide with the post-treatment period --- causal attribution is uncertain
    \vspace{0.4em}
    \item \textbf{Cannot separate} content preference shifts from changes in how the algorithm surfaces content
    \vspace{0.4em}
    \item \textbf{Trending $\neq$ feed} --- Twitter's trending list and the algorithmic feed are distinct systems; results speak to curation of ``what's trending,'' not all content delivery
\end{itemize}
\end{frame}

% ── SLIDE 13 — Key Takeaways ──────────────────────────────────────────────────
\begin{frame}{Key Takeaways}
\begin{enumerate}
    \setlength\itemsep{0.8em}
    \item \textbf{Twitter's algorithmic curation shifted significantly after the Musk acquisition} --- content category shares diverged from Reddit's organic baseline immediately after Oct 27, 2022
    \vspace{0.2em}
    \item \textbf{Sports, wrestling, combat sports, and political content were amplified;} female-coded and fandom content were suppressed --- a pattern consistent with platform demographic repositioning
    \vspace{0.2em}
    \item \textbf{Reddit organic activity serves as a viable counterfactual} for platform-level DiD, providing a clean pre-trend parallel and interpretable treatment effect estimates
\end{enumerate}
\end{frame}

% ── SLIDE 14 — Thank You ──────────────────────────────────────────────────────
\begin{frame}
\begin{center}
\vspace{1em}
{\Huge \textbf{Thank You}}\\[1em]
{\large Questions?}\\[2em]

\textbf{Alex Gamez}\\
\texttt{aGamez@utep.edu}\\[1em]

\textcolor{gray}{\small Data \& code available upon request}
\end{center}
\end{frame}

\end{document}
